%%%%%%%%%%%%%%%%%%%%%%%%%%%%%%%%%%%%%%%%%%%%%%%%%%%%%%%%%%%%%%%%%%%%%%%%%%%%%%%%
%%  EINLAGE 4: EINES IHRER ERGEBNISSE     Zwischenabgabe: 20.01.2027 (obligatorisch)
%%  Struktur: Beschreibung (1/4–1/2 S.) + Diskussion der Physik (1–3 S.) + Reflexion (1–2 S.)
%%  Rohdaten, Auswerteformeln, ausgewertete Daten, Grafiken gehören zur Einlage.
%%  Lange Messwertlisten -> Anhang (4_messergebnis_anhang.tex), nicht in den Text.
%%%%%%%%%%%%%%%%%%%%%%%%%%%%%%%%%%%%%%%%%%%%%%%%%%%%%%%%%%%%%%%%%%%%%%%%%%%%%%%%
\section{Einlage: Eines Ihrer Ergebnisse (Messung oder Berechnung).}
\label{sec:messergebnis}

\begin{hinweise}
Hier nehmen Sie ein von Ihnen gemessenes Resultat (Rohdaten, Auswerteformeln und die ausgewerteten Daten und
Grafiken) und diskutieren es. Ein Bezug zu Ihrem Paper (Einlage~2) ist Voraussetzung.
\end{hinweise}

\subsection{Beschreibung}
\begin{hinweise}
Die Beschreibung gibt alle notwendigen Informationen an, die die Einlage beschreiben, also zum Beispiel das
Datum, die Art der Messung und eine kurze Inhaltsübersicht, aber keine Nacherzählung (min. $1/4$ Seite bis
max. $1/2$ Seite!).
\end{hinweise}
\begin{umfang}{Einlage 4 (Messergebnis): Beschreibung}{0.25}{0.5}
\platzhalter{Datum, Art der Messung, Messaufbau in Kürze -- laut Arbeitsplan: 3D-gedruckter Wackelstein
mit eingebautem Inertialsensor (Gyroskop) auf nivellierter Glasplatte (30\,cm\,$\times$\,30\,cm),
reproduzierbarer Anstoß mit definiertem Anfangs-$\omega$; eine der Messreihen 1--3 (Anfangs-$\omega$
variiert, Nick- vs.\ Roll-Anregung, $\delta$ variiert); Auswertung mit der eigenen Python-Datenpipeline;
Rohdaten in Anhang~\ref{anh:messdaten}, Ergebnis in Abbildung~\ref{fig:messergebnis}}
\end{umfang}

\subsection{Diskussion der Physik}
\begin{hinweise}
Beschreiben Sie hier die Physik dieser Einlage.
\begin{itemize}\itemsep0pt
\item Was haben Sie gemessen oder berechnet?
\item Was ist die Hauptaussage?
\item Was sind die Voraussetzungen für Ihre Messung?
\item Welches sind Ihre Auswerteformeln?
\item Was trägt ihre \textbf{wissenschaftliche Arbeit} zur Argumentation bei?
\item Welche Physik dieser Einlage ist bei Ihrem Schwerpunktthema anwendbar?
\item Wo gab es Schwierigkeiten beim Verständnis der Physik? Welche?
\item Wie haben Sie Ihr Paper aus der Einlage \flqq Eine wissenschaftliche Arbeit\frqq{} für diese Einlage verwendet? Warum? Wo?
\item Stimmen Ihre Ergebnisse mit den Aussagen aus dem Paper überein?
\end{itemize}
\textbf{Die Diskussion der Physik sollte mindestens eine Seite und maximal 3 Seiten lang sein. Die
Längenangaben sind ohne Titel, Bilder, Tabellen und Tabellen- oder Bildlegenden und ohne übergrosse Abstände
zu Gleichungen zu verstehen. Ich werde gegebenenfalls Ihre Auswertung anhand der beigelegten Daten
überprüfen.}

\smallskip Regeln für Abbildungen/Tabellen/Gleichungen (Limmer): jede Abbildung mit Bildunterschrift und im
Text referenziert; Schrift in Diagrammen $\geq 10$\,pt in Endgröße; Achsen \enquote{Größe / Einheit}
(z.\,B. $\omega_z$\,/\,rad\,s$^{-1}$, nicht $\omega_z$ [rad/s]); bei guter Übereinstimmung
Messung--Theorie/Simulation zusätzlich Residuenplot; Tabellen mit Überschrift; Messdaten nicht (nur) als
Tabelle, sondern (auch) als Diagramm; wichtige Gleichungen nummeriert, alle Symbole definiert.
Die Gerüste unten (Gleichung, Abbildung, Tabelle) sind auf das Projekt zugeschnitten -- anpassen oder löschen.
Bezug zum Paper ist Pflicht: Messergebnis mit der Umkehrzeit-Formel bzw.\ den Simulationen aus
\cite{\PaperKey} und mit der eigenen Simulation (Arbeitspaket 1.6) vergleichen.
\end{hinweise}
\begin{umfang}{Einlage 4 (Messergebnis): Diskussion der Physik}{1}{3}
\platzhalter{Was wurde gemessen, Hauptaussage, Voraussetzungen (Rollen ohne Gleiten auf der Glasplatte,
Kalibrierung des Gyroskops, \dots)}

% ---------- Gerüst: Auswerteformel mit Nummer und Verweis ----------
% Wachstum der Nick- bzw. Rollschwingung in der instabilen Drehrichtung (Arbeitspaket 5.2):
\begin{equation}
  A(t) = A_0\,\mathrm{e}^{\gamma t},
  \label{eq:wachstum}
\end{equation}
wobei $A$ die Amplitude der Nick- bzw.\ Rollschwingung, $A_0$ ihr Anfangswert und $\gamma$ die
Wachstumsrate der Instabilität ist. \platzhalter{Gleichung prüfen/ersetzen; Umkehrzeit $t_r$ aus den
Gyroskopdaten definieren (z.\,B.\ Zeit bis zum Vorzeichenwechsel von $\omega_z$); Bezug zum Paper:
\enquote{Nach \cite[Gl.~(\dots)]{\PaperKey} gilt für $t_r$ \dots}; Gl.~\eqref{eq:wachstum} im Text verwenden}

% ---------- Gerüst: Abbildung (Datei nach abbildungen/ legen) ----------
\begin{figure}[htbp]
  \centering
  %\includegraphics[width=0.8\textwidth]{messergebnis.pdf}
  \fbox{\parbox[c][4cm][c]{0.7\textwidth}{\centering\platzhalter{Abbildung: abbildungen/messergebnis.pdf,
    z.\,B.\ $\omega_z(t)$ aus dem Gyroskop mit Spin-Umkehr}}}
  \caption{\platzhalter{Aussagekräftige Bildunterschrift: Winkelgeschwindigkeit $\omega_z$ um die
    Hochachse über der Zeit $t$ für \dots, Messbedingungen, Ergebnis der Anpassung mit
    Gl.~\eqref{eq:wachstum}. Fremde Abbildungen: \cite[Abb.~\dots]{\PaperKey}}}
  \label{fig:messergebnis}
\end{figure}

% ---------- Gerüst: kurze Ergebnistabelle (Rohdaten in den Anhang!) ----------
\begin{table}[htbp]
  \centering
  \caption{\platzhalter{Tabellenüberschrift, z.\,B.\ gemessene Umkehrzeit $t_r$ in Abhängigkeit
    vom Anfangs-$\omega_0$ (Messreihe 1) im Vergleich zu Simulation und Paper}}
  \label{tab:ergebnis}
  \begin{tabular}{cccc}
    \toprule
    {$\omega_0$ / \si{\radian\per\second}} & {$t_r$ (Messung) / \si{\second}}
      & {$t_r$ (Simulation) / \si{\second}} & {$t_r$ (Paper) / \si{\second}} \\
    \midrule
    \dots & \dots & \dots & \dots \\
    \bottomrule
  \end{tabular}
\end{table}

\platzhalter{Diskussion: Abbildung~\ref{fig:messergebnis} und Tabelle~\ref{tab:ergebnis} im Text
verwenden, Unsicherheiten (Fehlerrechnung, Arbeitspaket 5.3), Vergleich Messung -- eigene Simulation --
Paper (5.4), Anwendbarkeit auf \Schwerpunkt}
\end{umfang}

\subsection{Reflexion}
\begin{hinweise}
In der Reflexion sollen Sie sich Gedanken machen über die Wirkung dieser Einlage und der Arbeit mit ihr auf
Ihren Lernfortschritt.
\begin{itemize}\itemsep0pt
\item In welchem Zusammenhang steht die Einlage mit Ihrem Schwerpunktthema?
\item Was haben Sie bei der Erstellung dieser Einlage gelernt?
\item Wie hat Ihnen die Einlage beim Verständnis Ihres Projektthemas geholfen?
\item Wo gab es Schwierigkeiten bei der Beschäftigung mit der Einlage?
\item Was hätten Sie anders gemacht, wenn Sie die Einlage nochmals erstellen müssten?
\item Was sagen Ihre Kommilitoninnen oder Kommilitonen dazu, als Kommentar oder wenn Sie mit Ihnen über die Einlage diskutieren?
\item Haben Sie in der Durchführung oder Auswertung der Messung logische oder physikalische Fehler entdeckt?
\item Gab es bei der Beschäftigung mit der Einlage eine überraschende neue Erkenntnis?
\end{itemize}
\textbf{Die Reflexion für eine Einlage sollte mindestens eine Seite und maximal 2 Seiten lang sein. Die
Längenangaben sind ohne Titel, Bilder, Tabellen und Tabellen- oder Bildlegenden zu verstehen.}
\end{hinweise}
\begin{umfang}{Einlage 4 (Messergebnis): Reflexion}{1}{2}
\platzhalter{Reflexion in Fließtext -- Leitfragen erkennbar beantworten}
\end{umfang}

%\subsection{Eigene Bewertungskriterien}
%...
